% Paste after Section 4 of the interim report; do not add a document environment.
% Required packages in the main document: booktabs, tabularx, array, hyperref.
% Compile with XeLaTeX or LuaLaTeX and use a Vietnamese-capable font for names.
% If previewing these sections alone, set the section counter to 4 in the wrapper.
% Evidence checkpoint: main, dbbe536, 30 October 2026. Tests are historical records,
% not new executions performed when drafting this text on 1 November 2026.
% Sources: docs/validation/MVP-checks-2026-10-30.md; docs/deployment.md;
% frontend/docs/interim-student-checks.md; Git history; weekly progress Google Doc,
% tabs Weeks 39--40 and Weeks 41--43, read on 1 November 2026.
% The remaining plan is a proposed allocation based on recorded responsibilities.
% Confirm owners before submission. Add a prototype URL only after verifying it.

\section{Software Engineering, Quality and Trial Deployment}
\label{sec:quality}

\subsection{Source Control and Build Workflow}
Development is maintained in the
\href{https://github.com/huyho44/IELTS-Evaluator-Pro}{project repository}.
The interim implementation checkpoint is \texttt{main}, commit
\texttt{dbbe536}, dated 30 October 2026. Its history records database,
backend, learner-interface and instructor-review changes, including branch
integration. Package scripts provide repeatable lint, build and test commands.
These checks currently run locally; no repository CI workflow is present at
this checkpoint. Automated CI and more consistent issue-to-commit traceability
remain improvement tasks.

\subsection{Testing Strategy, Results and Known Defects}
Testing combines business-rule checks, isolated PostgreSQL integration tests,
frontend regression scripts and local API smoke tests. The
\href{https://github.com/huyho44/IELTS-Evaluator-Pro/blob/dbbe536/docs/validation/MVP-checks-2026-10-30.md}{30 October validation record}
reports the following results on Node.js~24.21.0.

\begin{center}
\begingroup
\small
\begin{tabularx}{\linewidth}{@{}p{0.30\linewidth}X@{}}
\toprule
\textbf{Check} & \textbf{Recorded result} \\
\midrule
Backend regular suite & 133 passed; 19 opt-in database tests skipped in this command. \\
PostgreSQL integration & All 19 passed in a separate isolated-database run. \\
Frontend regressions & All five scripts passed: session, catalogue, instructor, student learning and course completion. \\
Local API smoke tests & 11 endpoint checks passed across three roles, including rejection of learner access to operational reports. \\
Static checks and builds & Backend type checking, backend/frontend lint and production builds passed. \\
Deployment configuration & Docker Compose configuration validation passed. \\
\bottomrule
\end{tabularx}
\endgroup
\end{center}

The integration runner created a disposable database, applied migrations and
seed data, and removed that database afterwards. Scenarios covered repeated
submissions, stale drafts and review claims, publication timing, appeals,
invalid evidence, provider retries and worker recovery. A separate 28 October
browser check exercised selected learner flows and desktop/mobile layouts.
However, full browser acceptance of the integrated Writing Clinic across all
three roles remains pending; the existing root-endpoint test does not
establish that complete flow. A node-postgres warning about concurrent queries
on a transaction connection also requires investigation despite passing tests.
Mocked provider tests and synthetic evaluation replay verify software behaviour,
not AI scoring accuracy or learning improvement.

\subsection{Security, Privacy and Error Handling}
Passwords use salted \texttt{scrypt} hashes. Access tokens expire after
15 minutes; refresh tokens are stored as hashes and rotated on refresh.
Backend checks enforce roles, course assignments, enrolment and submission
ownership. Input and assessment validators reject invalid values and
unsupported quotations. Version checks and database-backed job recovery
protect against stale updates and repeated processing. AI failures are
recorded rather than replaced with fabricated scores; official Writing
results require teacher decisions. Provider credentials remain server-side,
and local environment files are excluded through ignore rules. Real essays
and labels must be collected with permission and excluded from Git; approved
learning materials and explicit disclosure of AI services are also required.
These measures do not constitute a completed security audit.

\subsection{Trial Deployment and Reproducibility}
The documented setup uses Node.js~24, PostgreSQL, migrations and an idempotent
demo seed. Deployment configuration separates the API, worker, frontend and
database; PostgreSQL is not exposed to the host, and Caddy provides the
planned HTTPS entry point. Operational reporting includes grading states,
provider errors, latency and token usage where available. However, the
validation environment had no running Docker daemon: image execution,
certificate issuance and public deployment were not verified. A reachable
prototype URL and secure trial-account instructions therefore remain
outstanding. Before final delivery, the team will verify a clean installation,
the three-role demo and backup/restore procedures in the deployment environment.

\section{Team Contributions and Remaining Plan}
\label{sec:plan}

\subsection{Individual Contributions and Evidence}
The
\href{https://docs.google.com/document/d/1tDeQUFaP0ttOuAwh-ihx1wdaSiAIPZxqVrUM3X8RfFQ/edit}{weekly progress log}
records the following division of work. Completed design-stage tasks and
recorded implementation responsibilities are distinguished below; an assigned
task alone is not treated as evidence of completion.

\begin{center}
\begingroup
\small
\begin{tabularx}{\linewidth}{@{}p{0.19\linewidth}p{0.105\linewidth}X@{}}
\toprule
\textbf{Member} & \textbf{ID} & \textbf{Recorded contribution or responsibility} \\
\midrule
Hồ Quốc Huy & 2352379 & Database setup with Hoàng Hải, marked complete in weeks 39--40; later assigned LLM integration research and code review. Git also records schema, review and MVP integration work (\texttt{98c99bc}, \texttt{dfac810}, \texttt{dbbe536}). \\
Phạm Quang Huy & 2352405 & UI/UX prototype with Hoàn Hải, marked complete in weeks 39--40; assigned frontend implementation and API integration in weeks 41--43. \\
Nguyễn Gia Hùng & 2352422 & Initial NestJS/database APIs, marked complete with \href{https://github.com/huyho44/IELTS-Evaluator-Pro/pull/12}{PR~12}; later assigned backend APIs with Hoàng Hải. \\
Đào Nguyên Hưng & 2452421 & Assigned React Router configuration, API-client setup and frontend-to-backend connection checks; completion evidence requires reconciliation. \\
Lê Đoàn Hoàn Hải & 2452307 & UI/UX prototype marked complete; learner catalogue, lessons, practice and session integration are also recorded in Git (\texttt{33b27a2}, \texttt{451c3be}, \texttt{2c31830}). \\
Phạm Hoàng Hải & 2450007 & Database setup with Quốc Huy, marked complete in weeks 39--40; assigned backend APIs with Gia Hùng in weeks 41--43. \\
\bottomrule
\end{tabularx}
\endgroup
\end{center}

The log links early deliverables to issues and pull requests, but several
later status and evidence fields remain incomplete. The team will reconcile
these fields with commits, test records and demonstrated outcomes instead
of assigning unsupported completion percentages.

\subsection{Risks and Recovery Actions}
The principal remaining risks are:
\begin{itemize}
  \item \textbf{Core-scope and integration gaps:} audit every mandatory LMS
  requirement and complete missing items before optional extensions.
  The frontend and backend groups will jointly verify the full three-role flow.
  \item \textbf{Deployment readiness:} resolve Docker availability and verify
  public access, HTTPS, migrations and trial accounts before final rehearsal.
  Quốc Huy will coordinate the proposed deployment work with the team.
  \item \textbf{Evaluation validity:} obtain permitted essays and independent
  expert labels, freeze development/held-out groups, and compare against
  declared baselines. Quốc Huy and Gia Hùng are proposed coordinators;
  synthetic fixtures will remain clearly separated from accuracy evidence.
\end{itemize}

\subsection{Remaining Milestones: Weeks 45--50}
The following plan proposes owners according to the recorded division of work.
Each milestone has a concrete acceptance outcome; individual allocations will
be confirmed in the weekly log.

\begin{center}
\begingroup
\small
\begin{tabularx}{\linewidth}{@{}p{0.10\linewidth}Xp{0.25\linewidth}@{}}
\toprule
\textbf{Week} & \textbf{Deliverable and acceptance outcome} & \textbf{Proposed owners} \\
\midrule
45--46 & Close mandatory-feature gaps and finish assessment/recommendation integration; complete a three-role scenario and prepare permitted evaluation data. & UI: Quang Huy, Hoàn Hải, Nguyên Hưng; APIs: Gia Hùng, Hoàng Hải; integration: Quốc Huy. \\
47 & Execute functional, security, usability and reliability checks; investigate the database warning and resolve critical defects with regression evidence. & All members, by module ownership. \\
48 & Verify public deployment and repeatable setup; run the declared evaluations and consolidate results, report, slides and demo script. & Quốc Huy coordinates; backend and frontend groups verify their flows. \\
49 & Rehearse the complete demo; verify lecturer access, sample data and trial accounts; resolve remaining blocking defects. & All six members. \\
50 & Deliver the final report, application, source, run instructions and evaluation/test evidence by 23:59, 13 December 2026 (UTC+7). & All six members; team leader consolidates submission. \\
\bottomrule
\end{tabularx}
\endgroup
\end{center}

\subsection{Interim Summary}
The interim checkpoint provides an integrated local implementation supported
by recorded automated checks. The remaining priorities are complete
three-role acceptance, verifiable public deployment, independent evaluation
of the advanced component and consolidated contribution evidence. These
outcomes, rather than the existence of screens or synthetic results alone,
will determine final readiness.
